\documentclass[a4paper, 11pt]{article}
\usepackage{amsmath, amsthm, breakurl, booktabs, colortbl, graphicx, listings, subcaption, tikz, xcolor}
\usepackage[inline]{enumitem}

%\usepackage[bitstream-charter]{mathdesign}
%\usepackage[T1]{fontenc}

\usepackage[colorlinks = none, hidelinks]{hyperref}

\DeclareCaptionFormat{custom}
{%
	\small{\textbf{\sffamily #1#2}\textit{ #3}}
}
\captionsetup{format=custom}

\allowdisplaybreaks

%%%%%%%%%%%%%%%%%
%% Definitions %%
%%%%%%%%%%%%%%%%%

% Code listings

\definecolor{codegreen}{rgb}{0,0.6,0}
\definecolor{codegray}{rgb}{0.5,0.5,0.5}
\definecolor{codepurple}{rgb}{0.58,0,0.82}
\definecolor{backcolour}{rgb}{0.95,0.95,0.92}

\lstdefinestyle{mystyle}{
	backgroundcolor=\color{backcolour},   
	commentstyle=\color{codegreen},
	keywordstyle=\color{magenta},
	numberstyle=\tiny\color{codegray},
	stringstyle=\color{codepurple},
	basicstyle=\ttfamily\footnotesize,
	breakatwhitespace=false,         
	breaklines=true,                 
	captionpos=b,                    
	keepspaces=true,                 
	numbers=left,                    
	numbersep=5pt,                  
	showspaces=false,                
	showstringspaces=false,
	showtabs=false,                  
	tabsize=2
}

\lstset{style=mystyle}

\definecolor{c1}{HTML}{faa08f}
\definecolor{c2}{HTML}{adadfe}

\theoremstyle{definition}
\newtheorem{exercise}{Exercise}

\title{Object--Oriented Python}
\author{Vassilis Markos}
\date{\today}

\begin{document}
	\maketitle
	
	\begin{enumerate}[label = E\arabic*.]
		\item Consider the following two python classes:
		\begin{lstlisting}[language=Python]
class A:
	def __init__(self, n):
		self.n = n

	def __eq__(self, __other):
		if not isinstance(__other, A):
			return False
		return self.n == __other.n

	class B(A):
		def __init__(self, n):
			self.n = n
		\end{lstlisting}
		\begin{enumerate}
			\item Explain what the above code snippet defines. How are the two classes defined? Is there any sort of inheritance? Are we overwriting any default python object methods?
			\item What will the following python script print?
			\begin{lstlisting}[language=Python]
if __name__ == "__main__":
	x = A(1)
	y = A(2)
	print(f"Are x and y equal? {x == y}")
			\end{lstlisting}
			\item What about the following?
			\begin{lstlisting}[language=Python]
if __name__ == "__main__":
	x = A(1)
	y = A(1)
	print(f"Are x and y equal? {x == y}")
			\end{lstlisting}
			\item What about the following?
			\begin{lstlisting}[language=Python]
if __name__ == "__main__":
	u = B(1)
	v = B(1)
	print(f"Are u and v equal? {u == v}")
			\end{lstlisting}
			\item What about the following?
			\begin{lstlisting}[language=Python]
if __name__ == "__main__":
	u = B(1)
	v = A(1)
	print(f"Are u and v equal? {u == v}")
			\end{lstlisting}
			\item What if we define the two classes as follows? How do the above cases change, if at all?
			\begin{lstlisting}[language=Python]
class A:
	def __init__(self, n):
		self.n = n

class B(A):
	def __init__(self, n):
		self.n = n

	def __eq__(self, __other):
		if not isinstance(__other, A):
			return False
		return self.n == __other.n
			\end{lstlisting}
		\end{enumerate}
		\item In this exercise we will create a (really--really) simple movie database management system.
		\begin{enumerate}
			\item As a first step, create a \texttt{Movie} class, with the following fields:
			\begin{itemize}
				\item \texttt{name}, a string representing the name of the movie;
				\item \texttt{genres}, a list of strings indicating the genre(s) of the movie;
				\item \texttt{duration}, an integer indicating the duration of the movie in minutes;
				\item \texttt{director}, a string containing the name of the movie director;
				\item \texttt{cast}, a list of strings, containing the names of the movie cast.
			\end{itemize}
			Moreover, the class should overwrite python's \texttt{\_\_str\_\_} method so as to print a movie as follows (with tabs, spacing and everything as shown below):
			\begin{verbatim}
				Movie: A Christmas Carol
				    Director: Robert Zemeckis
				    Genre   : Animation, Adventure, Comedy
				    Duration: 126 minutes
				    Cast    : Jim Carrey, Garry Oldman, Colin Firth
			\end{verbatim}
			You can read more about python's \texttt{\_\_str\_\_} magic function here: \href{https://docs.python.org/3/reference/datamodel.html#object.__str__}{\texttt{https://docs.python.org/3/reference/datamodel.html\#object.\_\_str\_\_}}
			\item Add an object method, \texttt{write\_to\_file()}, to class \texttt{Movie} that accepts a filename as an argument and adds a line to that file with all movie information in the following form:
			\begin{verbatim}
				name | director | genre | duration | cast
			\end{verbatim}
			where \texttt{|} is used as a delimiter, to distinguish between the various fields.
			For instance, for the movie ``A Christmas Carol'' shown above, this line should be as follows (split here just to fin in the page):
			\begin{verbatim}
				A Christmas Carol | Rober Zemeckis | Animation,
				Adventure, Comedy | 126 | Jim Carrey, Garry Oldman,
				Colin Firth
			\end{verbatim}
			\item Add an object method, \texttt{from\_text()}, that accepts a string of the following form:
			\begin{verbatim}
				name | director | genre | duration | cast
			\end{verbatim}
			and updates all class fields accordingly.
			
			\textit{If you wish, you can first read about how this can be achieved by using class methods\footnote{One related source: \href{https://stackoverflow.com/questions/5249813/create-python-objects-from-strings}{https://stackoverflow.com/questions/5249813/create-python-objects-from-strings}.} but it is not compulsory for this part of the exercise to use python's \texttt{@classmethod}.}
			\item Write a python script that:
			\begin{itemize}
				\item Asks the user to provide the name of a movie and then asks for additional information:
					\begin{itemize}
						\item movie director,
						\item genres,
						\item duration,
						\item cast.
					\end{itemize}
				\item It creates a \texttt{Movie} object with the provided values and writes it to a file (the file has to be created for the first time this script is used).
				\item It repeats the above until the user enters ``:q'' as the name of the movie.
			\end{itemize}
			\item Consider extending the above python script by adding several functionalities. These could include (but should not be restricted to) the following:
			\begin{itemize}
				\item duplicate checks, i.e. check whether a movie has already been written into the file;
				\item querying a movie, i.e., allowing the user not only to enter a movie but also ask for information about a movie;
				\item \ldots
			\end{itemize}
			Try to provide well--thought implementations. For instance, if you implement a querying functionality, what if the movie is not included in the file? How are you going to actually search the file?
		\end{enumerate}
	\end{enumerate}
\end{document}